%% notes-tex/database/spell-publications/sections/3-no-geo.tex
%% [[experiments.015-spell.publications]]
%% https://github.com/Mjvolk3/torchcell/tree/main/notes-tex/database/spell-publications/sections/3-no-geo.tex
%% SOURCE: numbers and table from
%% experiments/015-spell/scripts/spell_publication_tables.py; prose hand-authored

\section{Studies with no GEO accession\secstatus{todo}}\label{sec:no-geo}

\spellNNoGeo{} READMEs read \texttt{GEO ID: N/A} (Table~\ref{tab:no-geo}). These
studies were published between \spellNoGeoYearMin{} and \spellNoGeoYearMax{} and
account for \spellNNoGeoDatasets{} PCL files and \spellNNoGeoConditions{}
conditions. None of their PCL filenames contains a GEO series id, so the archive
gives no pointer to the authors' released data. The data location has to be
recovered from the paper: its methods, its supplementary files, or a laboratory
or repository address it names. \spellNNoGeoPmc{} of these studies have a PMC id
and \spellNNoGeoNoPmc{} do not. \spellNNoGeoTwoChannel{} of the \spellNNoGeo{}
are two-channel microarray studies. \spellNKoStudiesNoGeo{} carry a deletion
call, and \spellNNamedStudiesNoGeo{} of the \spellNNamedStudies{} gene-named
studies of Table~\ref{tab:gene-named}, Hughes 2000 and Mnaimneh 2004 among them,
are in this group.

All \spellNNoGeoPclPrefixed{} of their PCL filenames begin with \texttt{2010.}
and carry suffixes such as \texttt{flt}, \texttt{knn}, \texttt{avg},
\texttt{div} and \texttt{log}, which the READMEs do not define. Hypothesis
(untested): the suffixes name processing steps applied after the authors'
release, such as filtering, nearest-neighbor imputation, replicate averaging,
division by a reference and log transformation. If so, the PCL values for these
studies are further from the source measurements than a GEO series would be.

%% GENERATED FILE -- do not hand-edit.
%% SOURCE: experiments/015-spell/scripts/spell_publication_tables.py
\begin{landscape}
\begingroup
\footnotesize
\setlength{\tabcolsep}{4pt}
\renewcommand{\arraystretch}{1.1}
\begin{longtable}{@{}r@{\hspace{5pt}} L{82mm} L{40mm} L{60mm} L{28mm} r r r@{}}
\caption[]{The studies whose README reads \texttt{GEO ID: N/A}, in the order of Table~\ref{tab:all}. \emph{\#} is the row number in Table~\ref{tab:all} and \emph{KO genes} the estimate of Sec.~\ref{sec:knockout}. \emph{PCL files in the archive} lists each matrix SPELL distributes for the study, which is the only data location the archive gives for these rows. \emph{SPELL tags} are the topic tags in the README dataset table. A row with a PMC link has a full text in PubMed Central to search for the data location; a row without one does not.}
\label{tab:no-geo}\\
\toprule
\textbf{\#} & \textbf{Publication} & \textbf{Links} & \textbf{PCL files in the archive} & \textbf{SPELL tags} & \textbf{cond.} & \textbf{KO genes} & \textbf{ch.} \\
\midrule
\endfirsthead
\multicolumn{8}{@{}l}{\footnotesize\emph{Table~\ref{tab:no-geo}, continued}}\\
\toprule
\textbf{\#} & \textbf{Publication} & \textbf{Links} & \textbf{PCL files in the archive} & \textbf{SPELL tags} & \textbf{cond.} & \textbf{KO genes} & \textbf{ch.} \\
\midrule
\endhead
\bottomrule
\endfoot
1 & \textbf{Mnaimneh 2004}\newline \sourcetext{Mnaimneh S, Davierwala AP, Haynes J, et al., Hughes TR. Exploration of essential gene functions via titratable promoter alleles. \emph{Cell}} 2004;118(1):31--44. & \href{https://doi.org/10.1016/j.cell.2004.06.013}{doi:10.\allowbreak 1016/\allowbreak j.\allowbreak cell.\allowbreak 2004.\allowbreak 06.\allowbreak 013}\newline \href{https://pubmed.ncbi.nlm.nih.gov/15242642/}{PMID 15242642} & {\ttfamily\scriptsize 2010.\allowbreak EssentialGenes.\allowbreak pcl} & chemical stimulus; stress & 215 & -- & 2 \\
\addlinespace[3pt]
2 & \textbf{Hughes 2000}\newline \sourcetext{Hughes TR, Marton MJ, Jones AR, et al., Friend SH. Functional discovery via a compendium of expression profiles. \emph{Cell}} 2000;102(1):109--26. & \href{https://doi.org/10.1016/s0092-8674(00)00015-5}{doi:10.\allowbreak 1016/\allowbreak s0092-\allowbreak 8674(00)00015-\allowbreak 5}\newline \href{https://pubmed.ncbi.nlm.nih.gov/10929718/}{PMID 10929718} & {\ttfamily\scriptsize 2010.\allowbreak Hughes00.\allowbreak flt.\allowbreak knn.\allowbreak avg.\allowbreak pcl} & chemical stimulus & 300 & -- & 2 \\
\addlinespace[3pt]
20 & \textbf{Belli 2004}\newline \sourcetext{Bell\'{\i} G, Molina MM, Garc\'{\i}a-Mart\'{\i}nez J, P\'erez-Ort\'{\i}n JE, Herrero E. Saccharomyces cerevisiae glutaredoxin 5-deficient cells subjected to continuous oxidizing conditions are affected in the expression of specific sets of genes. \emph{J Biol Chem}} 2004;279(13):12386--95. & \href{https://doi.org/10.1074/jbc.M311879200}{doi:10.\allowbreak 1074/\allowbreak jbc.\allowbreak M311879200}\newline \href{https://pubmed.ncbi.nlm.nih.gov/14722110/}{PMID 14722110} & {\ttfamily\scriptsize 2010.\allowbreak Belli04.\allowbreak filter.\allowbreak flt.\allowbreak knn.\allowbreak avg.\allowbreak div.\allowbreak log.\allowbreak pcl} & oxidative stress; stress & 9 & 2 & 2 \\
\addlinespace[3pt]
29 & \textbf{Huang 2002}\newline \sourcetext{Huang D, Moffat J, Andrews B. Dissection of a complex phenotype by functional genomics reveals roles for the yeast cyclin-dependent protein kinase Pho85 in stress adaptation and cell integrity. \emph{Mol Cell Biol}} 2002;22(14):5076--88. & \href{https://doi.org/10.1128/MCB.22.14.5076-5088.2002}{doi:10.\allowbreak 1128/\allowbreak MCB.\allowbreak 22.\allowbreak 14.\allowbreak 5076-\allowbreak 5088.\allowbreak 2002}\newline \href{https://pubmed.ncbi.nlm.nih.gov/12077337/}{PMID 12077337}\newline \href{https://pmc.ncbi.nlm.nih.gov/articles/PMC139770/}{PMC139770} & {\ttfamily\scriptsize 2010.\allowbreak Huang02.\allowbreak flt.\allowbreak knn.\allowbreak avg.\allowbreak pcl} & protein phosphorylation & 20 & 12 & 2 \\
\addlinespace[3pt]
31 & \textbf{Pitkanen 2004}\newline \sourcetext{Pitk\"anen JP, T\"orm\"a A, Alff S, Huopaniemi L, Mattila P, Renkonen R. Excess mannose limits the growth of phosphomannose isomerase PMI40 deletion strain of Saccharomyces cerevisiae. \emph{J Biol Chem}} 2004;279(53):55737--43. & \href{https://doi.org/10.1074/jbc.M410619200}{doi:10.\allowbreak 1074/\allowbreak jbc.\allowbreak M410619200}\newline \href{https://pubmed.ncbi.nlm.nih.gov/15520001/}{PMID 15520001} & {\ttfamily\scriptsize 2010.\allowbreak Pitkanen04.\allowbreak filter.\allowbreak flt.\allowbreak knn.\allowbreak avg.\allowbreak div.\allowbreak log.\allowbreak pcl} & carbon utilization; protein glycosylation & 15 & 1 & 2 \\
\addlinespace[3pt]
36 & \textbf{Lee 2005}\newline \sourcetext{Lee A, Henras AK, Chanfreau G. Multiple RNA surveillance pathways limit aberrant expression of iron uptake mRNAs and prevent iron toxicity in S. cerevisiae. \emph{Mol Cell}} 2005;19(1):39--51. & \href{https://doi.org/10.1016/j.molcel.2005.05.021}{doi:10.\allowbreak 1016/\allowbreak j.\allowbreak molcel.\allowbreak 2005.\allowbreak 05.\allowbreak 021}\newline \href{https://pubmed.ncbi.nlm.nih.gov/15989963/}{PMID 15989963} & {\ttfamily\scriptsize 2010.\allowbreak Lee05.\allowbreak filter.\allowbreak flt.\allowbreak knn.\allowbreak avg.\allowbreak div.\allowbreak log.\allowbreak pcl} & RNA catabolism & 9 & 1 & 2 \\
\addlinespace[3pt]
40 & \textbf{Eriksson 2005}\newline \sourcetext{Eriksson PR, Mendiratta G, McLaughlin NB, et al., Clark DJ. Global regulation by the yeast Spt10 protein is mediated through chromatin structure and the histone upstream activating sequence elements. \emph{Mol Cell Biol}} 2005;25(20):9127--37. & \href{https://doi.org/10.1128/MCB.25.20.9127-9137.2005}{doi:10.\allowbreak 1128/\allowbreak MCB.\allowbreak 25.\allowbreak 20.\allowbreak 9127-\allowbreak 9137.\allowbreak 2005}\newline \href{https://pubmed.ncbi.nlm.nih.gov/16199888/}{PMID 16199888}\newline \href{https://pmc.ncbi.nlm.nih.gov/articles/PMC1265784/}{PMC1265784} & {\ttfamily\scriptsize 2010.\allowbreak Eriksson05.\allowbreak filter.\allowbreak flt.\allowbreak knn.\allowbreak avg.\allowbreak div.\allowbreak log.\allowbreak pcl} & chromatin organization; histone modification & 6 & 1 & 2 \\
\addlinespace[3pt]
45 & \textbf{Rudra 2005}\newline \sourcetext{Rudra D, Zhao Y, Warner JR. Central role of Ifh1p-Fhl1p interaction in the synthesis of yeast ribosomal proteins. \emph{EMBO J}} 2005;24(3):533--42. & \href{https://doi.org/10.1038/sj.emboj.7600553}{doi:10.\allowbreak 1038/\allowbreak sj.\allowbreak emboj.\allowbreak 7600553}\newline \href{https://pubmed.ncbi.nlm.nih.gov/15692568/}{PMID 15692568}\newline \href{https://pmc.ncbi.nlm.nih.gov/articles/PMC548658/}{PMC548658} & {\ttfamily\scriptsize 2010.\allowbreak Rudra05.\allowbreak filter.\allowbreak flt.\allowbreak knn.\allowbreak avg.\allowbreak div.\allowbreak log.\allowbreak pcl} & transcription & 6 & 1 & 1 \\
\addlinespace[3pt]
63 & \textbf{Gasch 2000}\newline \sourcetext{Gasch AP, Spellman PT, Kao CM, et al., Brown PO. Genomic expression programs in the response of yeast cells to environmental changes. \emph{Mol Biol Cell}} 2000;11(12):4241--57. & \href{https://doi.org/10.1091/mbc.11.12.4241}{doi:10.\allowbreak 1091/\allowbreak mbc.\allowbreak 11.\allowbreak 12.\allowbreak 4241}\newline \href{https://pubmed.ncbi.nlm.nih.gov/11102521/}{PMID 11102521}\newline \href{https://pmc.ncbi.nlm.nih.gov/articles/PMC15070/}{PMC15070} & {\ttfamily\scriptsize 2010.\allowbreak Gasch00\_\allowbreak hyper-\allowbreak osmotic.\allowbreak flt.\allowbreak knn.\allowbreak avg.\allowbreak pcl\newline 2010.\allowbreak Gasch00\_\allowbreak hypo-\allowbreak osmotic.\allowbreak flt.\allowbreak knn.\allowbreak avg.\allowbreak pcl\newline 2010.\allowbreak Gasch00\_\allowbreak menadione.\allowbreak flt.\allowbreak knn.\allowbreak avg.\allowbreak pcl\newline 2010.\allowbreak Gasch00\_\allowbreak Ndepletion.\allowbreak flt.\allowbreak knn.\allowbreak avg.\allowbreak pcl\newline 2010.\allowbreak Gasch00\_\allowbreak stationaryPhase(y12).\allowbreak flt.\allowbreak knn.\allowbreak avg.\allowbreak pcl\newline 2010.\allowbreak Gasch00\_\allowbreak stationaryPhase(y14).\allowbreak flt.\allowbreak knn.\allowbreak avg.\allowbreak pcl\newline 2010.\allowbreak Gasch00\_\allowbreak steadyState(y13).\allowbreak flt.\allowbreak knn.\allowbreak avg.\allowbreak pcl\newline 2010.\allowbreak Gasch00\_\allowbreak steadyState(y14).\allowbreak flt.\allowbreak knn.\allowbreak avg.\allowbreak pcl\newline 2010.\allowbreak Gasch00\_\allowbreak HS29-\allowbreak 33.\allowbreak flt.\allowbreak knn.\allowbreak avg.\allowbreak pcl\newline 2010.\allowbreak Gasch00\_\allowbreak HS30-\allowbreak 37.\allowbreak flt.\allowbreak knn.\allowbreak avg.\allowbreak pcl\newline 2010.\allowbreak Gasch00\_\allowbreak HS37-\allowbreak 25.\allowbreak flt.\allowbreak knn.\allowbreak avg.\allowbreak pcl\newline 2010.\allowbreak Gasch00\_\allowbreak HSmild.\allowbreak flt.\allowbreak knn.\allowbreak avg.\allowbreak pcl\newline 2010.\allowbreak Gasch00\_\allowbreak HSto37.\allowbreak flt.\allowbreak knn.\allowbreak avg.\allowbreak pcl\newline 2010.\allowbreak Gasch00\_\allowbreak DTT(y13).\allowbreak flt.\allowbreak knn.\allowbreak avg.\allowbreak pcl\newline 2010.\allowbreak Gasch00\_\allowbreak DTT(y14).\allowbreak flt.\allowbreak knn.\allowbreak avg.\allowbreak pcl\newline 2010.\allowbreak Gasch00\_\allowbreak HOtimeCourse.\allowbreak flt.\allowbreak knn.\allowbreak avg.\allowbreak pcl\newline 2010.\allowbreak Gasch00\_\allowbreak HS25-\allowbreak 37.\allowbreak flt.\allowbreak knn.\allowbreak avg.\allowbreak pcl\newline 2010.\allowbreak Gasch00\_\allowbreak carbonSources.\allowbreak flt.\allowbreak knn.\allowbreak avg.\allowbreak pcl\newline 2010.\allowbreak Gasch00\_\allowbreak diamideTreatment.\allowbreak flt.\allowbreak knn.\allowbreak avg.\allowbreak pcl\newline 2010.\allowbreak Gasch00\_\allowbreak adenineStarvation.\allowbreak flt.\allowbreak knn.\allowbreak avg.\allowbreak pcl} & amino acid utilization; carbon utilization; chemical stimulus; heat shock; nitrogen utilization; osmotic stress; oxidative stress; response to unfolded protein; starvation; stationary phase entry; stress & 138 & -- & 2 \\
\addlinespace[3pt]
65 & \textbf{O'Rourke 2004}\newline \sourcetext{O'Rourke SM, Herskowitz I. Unique and redundant roles for HOG MAPK pathway components as revealed by whole-genome expression analysis. \emph{Mol Biol Cell}} 2004;15(2):532--42. & \href{https://doi.org/10.1091/mbc.e03-07-0521}{doi:10.\allowbreak 1091/\allowbreak mbc.\allowbreak e03-\allowbreak 07-\allowbreak 0521}\newline \href{https://pubmed.ncbi.nlm.nih.gov/14595107/}{PMID 14595107}\newline \href{https://pmc.ncbi.nlm.nih.gov/articles/PMC329229/}{PMC329229} & {\ttfamily\scriptsize 2010.\allowbreak ORourke03.\allowbreak flt.\allowbreak knn.\allowbreak avg.\allowbreak pcl} & osmotic stress; protein phosphorylation; signaling; stress & 133 & -- & 2 \\
\addlinespace[3pt]
66 & \textbf{Brem 2005}\newline \sourcetext{Brem RB, Kruglyak L. The landscape of genetic complexity across 5,700 gene expression traits in yeast. \emph{Proc Natl Acad Sci U S A}} 2005;102(5):1572--7. & \href{https://doi.org/10.1073/pnas.0408709102}{doi:10.\allowbreak 1073/\allowbreak pnas.\allowbreak 0408709102}\newline \href{https://pubmed.ncbi.nlm.nih.gov/15659551/}{PMID 15659551}\newline \href{https://pmc.ncbi.nlm.nih.gov/articles/PMC547855/}{PMC547855} & {\ttfamily\scriptsize 2010.\allowbreak Brem05\_\allowbreak orig.\allowbreak flt.\allowbreak knn.\allowbreak avg.\allowbreak pcl} & QTLs & 131 & -- & 2 \\
\addlinespace[3pt]
69 & \textbf{Saldanha 2004}\newline \sourcetext{Saldanha AJ, Brauer MJ, Botstein D. Nutritional homeostasis in batch and steady-state culture of yeast. \emph{Mol Biol Cell}} 2004;15(9):4089--104. & \href{https://doi.org/10.1091/mbc.e04-04-0306}{doi:10.\allowbreak 1091/\allowbreak mbc.\allowbreak e04-\allowbreak 04-\allowbreak 0306}\newline \href{https://pubmed.ncbi.nlm.nih.gov/15240820/}{PMID 15240820}\newline \href{https://pmc.ncbi.nlm.nih.gov/articles/PMC515343/}{PMC515343} & {\ttfamily\scriptsize 2010.\allowbreak Saldanha04\_\allowbreak UracilBatchChem.\allowbreak flt.\allowbreak knn.\allowbreak avg.\allowbreak pcl\newline 2010.\allowbreak Saldanha04\_\allowbreak UraSulPhoLeuComp.\allowbreak flt.\allowbreak knn.\allowbreak avg.\allowbreak pcl\newline 2010.\allowbreak Saldanha04\_\allowbreak SulfateBatchChem.\allowbreak flt.\allowbreak knn.\allowbreak avg.\allowbreak pcl\newline 2010.\allowbreak Saldanha04\_\allowbreak LeucineBatchChem.\allowbreak flt.\allowbreak knn.\allowbreak avg.\allowbreak pcl\newline 2010.\allowbreak Saldanha04\_\allowbreak PhosphateBatchChem.\allowbreak flt.\allowbreak knn.\allowbreak avg.\allowbreak pcl} & amino acid utilization; nutrient utilization; phosphorus utilization; sulfur utilization & 124 & -- & 2 \\
\addlinespace[3pt]
82 & \textbf{Shapira 2004}\newline \sourcetext{Shapira M, Segal E, Botstein D. Disruption of yeast forkhead-associated cell cycle transcription by oxidative stress. \emph{Mol Biol Cell}} 2004;15(12):5659--69. & \href{https://doi.org/10.1091/mbc.e04-04-0340}{doi:10.\allowbreak 1091/\allowbreak mbc.\allowbreak e04-\allowbreak 04-\allowbreak 0340}\newline \href{https://pubmed.ncbi.nlm.nih.gov/15371544/}{PMID 15371544}\newline \href{https://pmc.ncbi.nlm.nih.gov/articles/PMC532044/}{PMC532044} & {\ttfamily\scriptsize 2010.\allowbreak Shapira04.\allowbreak flt.\allowbreak knn.\allowbreak avg.\allowbreak pcl} & chemical stimulus; mitotic cell cycle; oxidative stress; stress & 70 & -- & 2 \\
\addlinespace[3pt]
95 & \textbf{Spellman 1998}\newline \sourcetext{Spellman PT, Sherlock G, Zhang MQ, et al., Futcher B. Comprehensive identification of cell cycle-regulated genes of the yeast Saccharomyces cerevisiae by microarray hybridization. \emph{Mol Biol Cell}} 1998;9(12):3273--97. & \href{https://doi.org/10.1091/mbc.9.12.3273}{doi:10.\allowbreak 1091/\allowbreak mbc.\allowbreak 9.\allowbreak 12.\allowbreak 3273}\newline \href{https://pubmed.ncbi.nlm.nih.gov/9843569/}{PMID 9843569}\newline \href{https://pmc.ncbi.nlm.nih.gov/articles/PMC25624/}{PMC25624} & {\ttfamily\scriptsize 2010.\allowbreak Spellman98\_\allowbreak cdc15.\allowbreak flt.\allowbreak knn.\allowbreak avg.\allowbreak pcl\newline 2010.\allowbreak Spellman98\_\allowbreak elutriation.\allowbreak flt.\allowbreak knn.\allowbreak avg.\allowbreak pcl\newline 2010.\allowbreak Spellman98\_\allowbreak alphaFactor.\allowbreak flt.\allowbreak knn.\allowbreak avg.\allowbreak pcl} & mitotic cell cycle & 55 & -- & 2 \\
\addlinespace[3pt]
97 & \textbf{Gasch 2001}\newline \sourcetext{Gasch AP, Huang M, Metzner S, Botstein D, Elledge SJ, Brown PO. Genomic expression responses to DNA-damaging agents and the regulatory role of the yeast ATR homolog Mec1p. \emph{Mol Biol Cell}} 2001;12(10):2987--3003. & \href{https://doi.org/10.1091/mbc.12.10.2987}{doi:10.\allowbreak 1091/\allowbreak mbc.\allowbreak 12.\allowbreak 10.\allowbreak 2987}\newline \href{https://pubmed.ncbi.nlm.nih.gov/11598186/}{PMID 11598186}\newline \href{https://pmc.ncbi.nlm.nih.gov/articles/PMC60150/}{PMC60150} & {\ttfamily\scriptsize 2010.\allowbreak DNAdamage.\allowbreak pcl} & DNA damage stimulus; radiation; signaling; stress & 52 & -- & 2 \\
\addlinespace[3pt]
100 & \textbf{Lai 2006}\newline \sourcetext{Lai LC, Kosorukoff AL, Burke PV, Kwast KE. Metabolic-state-dependent remodeling of the transcriptome in response to anoxia and subsequent reoxygenation in Saccharomyces cerevisiae. \emph{Eukaryot Cell}} 2006;5(9):1468--89. & \href{https://doi.org/10.1128/EC.00107-06}{doi:10.\allowbreak 1128/\allowbreak EC.\allowbreak 00107-\allowbreak 06}\newline \href{https://pubmed.ncbi.nlm.nih.gov/16963631/}{PMID 16963631}\newline \href{https://pmc.ncbi.nlm.nih.gov/articles/PMC1563586/}{PMC1563586} & {\ttfamily\scriptsize 2010.\allowbreak glucose\_\allowbreak anoxia\_\allowbreak reoxygenation.\allowbreak pcl\newline 2010.\allowbreak galactose\_\allowbreak anoxia\_\allowbreak reoxygenation.\allowbreak pcl} & carbon utilization; oxygen level alteration; stress & 48 & -- & 2 \\
\addlinespace[3pt]
105 & \textbf{Causton 2001}\newline \sourcetext{Causton HC, Ren B, Koh SS, et al., Young RA. Remodeling of yeast genome expression in response to environmental changes. \emph{Mol Biol Cell}} 2001;12(2):323--37. & \href{https://doi.org/10.1091/mbc.12.2.323}{doi:10.\allowbreak 1091/\allowbreak mbc.\allowbreak 12.\allowbreak 2.\allowbreak 323}\newline \href{https://pubmed.ncbi.nlm.nih.gov/11179418/}{PMID 11179418}\newline \href{https://pmc.ncbi.nlm.nih.gov/articles/PMC30946/}{PMC30946} & {\ttfamily\scriptsize 2010.\allowbreak Causton01\_\allowbreak peroxide.\allowbreak filter.\allowbreak flt.\allowbreak knn.\allowbreak avg.\allowbreak div.\allowbreak log.\allowbreak pcl\newline 2010.\allowbreak Causton01\_\allowbreak sorbitol.\allowbreak filter.\allowbreak flt.\allowbreak knn.\allowbreak avg.\allowbreak div.\allowbreak log.\allowbreak pcl\newline 2010.\allowbreak Causton01\_\allowbreak acid.\allowbreak filter.\allowbreak flt.\allowbreak knn.\allowbreak avg.\allowbreak div.\allowbreak log.\allowbreak pcl\newline 2010.\allowbreak Causton01\_\allowbreak alkali.\allowbreak filter.\allowbreak flt.\allowbreak knn.\allowbreak avg.\allowbreak div.\allowbreak log.\allowbreak pcl\newline 2010.\allowbreak Causton01\_\allowbreak heat.\allowbreak filter.\allowbreak flt.\allowbreak knn.\allowbreak avg.\allowbreak div.\allowbreak log.\allowbreak pcl\newline 2010.\allowbreak Causton01\_\allowbreak NaCl.\allowbreak filter.\allowbreak flt.\allowbreak knn.\allowbreak avg.\allowbreak div.\allowbreak log.\allowbreak pcl} & chemical stimulus; heat shock; osmotic stress; oxidative stress; response to unfolded protein; stress & 45 & -- & 2 \\
\addlinespace[3pt]
115 & \textbf{Yoshimoto 2002}\newline \sourcetext{Yoshimoto H, Saltsman K, Gasch AP, et al., Cyert MS. Genome-wide analysis of gene expression regulated by the calcineurin/Crz1p signaling pathway in Saccharomyces cerevisiae. \emph{J Biol Chem}} 2002;277(34):31079--88. & \href{https://doi.org/10.1074/jbc.M202718200}{doi:10.\allowbreak 1074/\allowbreak jbc.\allowbreak M202718200}\newline \href{https://pubmed.ncbi.nlm.nih.gov/12058033/}{PMID 12058033} & {\ttfamily\scriptsize 2010.\allowbreak Yoshimoto02\_\allowbreak Ca.\allowbreak flt.\allowbreak knn.\allowbreak avg.\allowbreak pcl\newline 2010.\allowbreak Yoshimoto02\_\allowbreak Na.\allowbreak flt.\allowbreak knn.\allowbreak avg.\allowbreak pcl} & cellular ion homeostasis; chemical stimulus; signaling & 40 & -- & 2 \\
\addlinespace[3pt]
137 & \textbf{Segal 2003}\newline \sourcetext{Segal E, Shapira M, Regev A, et al., Friedman N. Module networks: identifying regulatory modules and their condition-specific regulators from gene expression data. \emph{Nat Genet}} 2003;34(2):166--76. & \href{https://doi.org/10.1038/ng1165}{doi:10.\allowbreak 1038/\allowbreak ng1165}\newline \href{https://pubmed.ncbi.nlm.nih.gov/12740579/}{PMID 12740579} & {\ttfamily\scriptsize 2010.\allowbreak Segal03\_\allowbreak HSkin82.\allowbreak flt.\allowbreak knn.\allowbreak avg.\allowbreak pcl\newline 2010.\allowbreak Segal03\_\allowbreak hypo-\allowbreak osmoticPpt1.\allowbreak flt.\allowbreak knn.\allowbreak avg.\allowbreak pcl\newline 2010.\allowbreak Segal03\_\allowbreak stationaryPhaseYPL230W.\allowbreak flt.\allowbreak knn.\allowbreak avg.\allowbreak pcl} & heat shock; osmotic stress; protein dephosphorylation; protein phosphorylation; stationary phase entry; stress & 32 & -- & 2 \\
\addlinespace[3pt]
140 & \textbf{Duvel 2003}\newline \sourcetext{D\"uvel K, Santhanam A, Garrett S, Schneper L, Broach JR. Multiple roles of Tap42 in mediating rapamycin-induced transcriptional changes in yeast. \emph{Mol Cell}} 2003;11(6):1467--78. & \href{https://doi.org/10.1016/s1097-2765(03)00228-4}{doi:10.\allowbreak 1016/\allowbreak s1097-\allowbreak 2765(03)00228-\allowbreak 4}\newline \href{https://pubmed.ncbi.nlm.nih.gov/12820961/}{PMID 12820961} & {\ttfamily\scriptsize 2010.\allowbreak Duvel03\_\allowbreak delayedRap.\allowbreak flt.\allowbreak knn.\allowbreak avg.\allowbreak pcl\newline 2010.\allowbreak Duvel03\_\allowbreak immediateRap.\allowbreak flt.\allowbreak knn.\allowbreak avg.\allowbreak pcl} & chemical stimulus; heat shock; protein dephosphorylation; signaling & 30 & -- & 2 \\
\addlinespace[3pt]
141 & \textbf{Fleming 2002}\newline \sourcetext{Fleming JA, Lightcap ES, Sadis S, Thoroddsen V, Bulawa CE, Blackman RK. Complementary whole-genome technologies reveal the cellular response to proteasome inhibition by PS-341. \emph{Proc Natl Acad Sci U S A}} 2002;99(3):1461--6. & \href{https://doi.org/10.1073/pnas.032516399}{doi:10.\allowbreak 1073/\allowbreak pnas.\allowbreak 032516399}\newline \href{https://pubmed.ncbi.nlm.nih.gov/11830665/}{PMID 11830665}\newline \href{https://pmc.ncbi.nlm.nih.gov/articles/PMC122213/}{PMC122213} & {\ttfamily\scriptsize 2010.\allowbreak Fleming99.\allowbreak log.\allowbreak flt.\allowbreak knn.\allowbreak avg.\allowbreak pcl} & chemical stimulus; proteolysis & 30 & -- & 2 \\
\addlinespace[3pt]
165 & \textbf{Sapra 2004}\newline \sourcetext{Sapra AK, Arava Y, Khandelia P, Vijayraghavan U. Genome-wide analysis of pre-mRNA splicing: intron features govern the requirement for the second-step factor, Prp17 in Saccharomyces cerevisiae and Schizosaccharomyces pombe. \emph{J Biol Chem}} 2004;279(50):52437--46. & \href{https://doi.org/10.1074/jbc.M408815200}{doi:10.\allowbreak 1074/\allowbreak jbc.\allowbreak M408815200}\newline \href{https://pubmed.ncbi.nlm.nih.gov/15452114/}{PMID 15452114} & {\ttfamily\scriptsize 2010.\allowbreak Sapra04.\allowbreak flt.\allowbreak knn.\allowbreak avg.\allowbreak pcl} & mRNA processing & 24 & 1 & 2 \\
\addlinespace[3pt]
168 & \textbf{Caba 2005}\newline \sourcetext{Caba E, Dickinson DA, Warnes GR, Aubrecht J. Differentiating mechanisms of toxicity using global gene expression analysis in Saccharomyces cerevisiae. \emph{Mutat Res}} 2005;575(1--2):34--46. & \href{https://doi.org/10.1016/j.mrfmmm.2005.02.005}{doi:10.\allowbreak 1016/\allowbreak j.\allowbreak mrfmmm.\allowbreak 2005.\allowbreak 02.\allowbreak 005}\newline \href{https://pubmed.ncbi.nlm.nih.gov/15878181/}{PMID 15878181} & {\ttfamily\scriptsize 2010.\allowbreak Caba05.\allowbreak filter.\allowbreak flt.\allowbreak knn.\allowbreak avg.\allowbreak div.\allowbreak log.\allowbreak pcl} & DNA damage stimulus; osmotic stress; stress & 24 & -- & 1 \\
\addlinespace[3pt]
176 & \textbf{Meneghini 2003}\newline \sourcetext{Meneghini MD, Wu M, Madhani HD. Conserved histone variant H2A.Z protects euchromatin from the ectopic spread of silent heterochromatin. \emph{Cell}} 2003;112(5):725--36. & \href{https://doi.org/10.1016/s0092-8674(03)00123-5}{doi:10.\allowbreak 1016/\allowbreak s0092-\allowbreak 8674(03)00123-\allowbreak 5}\newline \href{https://pubmed.ncbi.nlm.nih.gov/12628191/}{PMID 12628191} & {\ttfamily\scriptsize 2010.\allowbreak swr1\_\allowbreak htz1\_\allowbreak hmr\_\allowbreak sir2.\allowbreak pcl} & chromatin organization & 24 & -- & 2 \\
\addlinespace[3pt]
177 & \textbf{Primig 2000}\newline \sourcetext{Primig M, Williams RM, Winzeler EA, et al., Esposito RE. The core meiotic transcriptome in budding yeasts. \emph{Nat Genet}} 2000;26(4):415--23. & \href{https://doi.org/10.1038/82539}{doi:10.\allowbreak 1038/\allowbreak 82539}\newline \href{https://pubmed.ncbi.nlm.nih.gov/11101837/}{PMID 11101837} & {\ttfamily\scriptsize 2010.\allowbreak Primig00.\allowbreak filter.\allowbreak flt.\allowbreak knn.\allowbreak avg.\allowbreak div.\allowbreak log.\allowbreak pcl} & sporulation & 24 & -- & 2 \\
\addlinespace[3pt]
178 & \textbf{Schawalder 2004}\newline \sourcetext{Schawalder SB, Kabani M, Howald I, Choudhury U, Werner M, Shore D. Growth-regulated recruitment of the essential yeast ribosomal protein gene activator Ifh1. \emph{Nature}} 2004;432(7020):1058--61. & \href{https://doi.org/10.1038/nature03200}{doi:10.\allowbreak 1038/\allowbreak nature03200}\newline \href{https://pubmed.ncbi.nlm.nih.gov/15616569/}{PMID 15616569} & {\ttfamily\scriptsize 2010.\allowbreak Schawalder04.\allowbreak filter.\allowbreak flt.\allowbreak knn.\allowbreak avg.\allowbreak div.\allowbreak log.\allowbreak pcl} & transcription & 24 & -- & 2 \\
\addlinespace[3pt]
180 & \textbf{Tai 2005}\newline \sourcetext{Tai SL, Boer VM, Daran-Lapujade P, et al., Pronk JT. Two-dimensional transcriptome analysis in chemostat cultures. Combinatorial effects of oxygen availability and macronutrient limitation in Saccharomyces cerevisiae. \emph{J Biol Chem}} 2005;280(1):437--47. & \href{https://doi.org/10.1074/jbc.M410573200}{doi:10.\allowbreak 1074/\allowbreak jbc.\allowbreak M410573200}\newline \href{https://pubmed.ncbi.nlm.nih.gov/15496405/}{PMID 15496405} & {\ttfamily\scriptsize 2010.\allowbreak Tai05.\allowbreak filter.\allowbreak flt.\allowbreak knn.\allowbreak avg.\allowbreak div.\allowbreak log.\allowbreak pcl} & carbon utilization; fermentation; nitrogen utilization; oxygen level alteration; phosphorus utilization; respiration; sulfur utilization & 24 & -- & 2 \\
\addlinespace[3pt]
194 & \textbf{Ideker 2001}\newline \sourcetext{Ideker T, Thorsson V, Ranish JA, et al., Hood L. Integrated genomic and proteomic analyses of a systematically perturbed metabolic network. \emph{Science}} 2001;292(5518):929--34. & \href{https://doi.org/10.1126/science.292.5518.929}{doi:10.\allowbreak 1126/\allowbreak science.\allowbreak 292.\allowbreak 5518.\allowbreak 929}\newline \href{https://pubmed.ncbi.nlm.nih.gov/11340206/}{PMID 11340206} & {\ttfamily\scriptsize 2010.\allowbreak Ideker01.\allowbreak flt.\allowbreak knn.\allowbreak avg.\allowbreak pcl} & carbon utilization & 21 & -- & 2 \\
\addlinespace[3pt]
198 & \textbf{Brauer 2005}\newline \sourcetext{Brauer MJ, Saldanha AJ, Dolinski K, Botstein D. Homeostatic adjustment and metabolic remodeling in glucose-limited yeast cultures. \emph{Mol Biol Cell}} 2005;16(5):2503--17. & \href{https://doi.org/10.1091/mbc.e04-11-0968}{doi:10.\allowbreak 1091/\allowbreak mbc.\allowbreak e04-\allowbreak 11-\allowbreak 0968}\newline \href{https://pubmed.ncbi.nlm.nih.gov/15758028/}{PMID 15758028}\newline \href{https://pmc.ncbi.nlm.nih.gov/articles/PMC1087253/}{PMC1087253} & {\ttfamily\scriptsize 2010.\allowbreak Brauer05\_\allowbreak batch1.\allowbreak flt.\allowbreak knn.\allowbreak avg.\allowbreak pcl\newline 2010.\allowbreak Brauer05\_\allowbreak batch2.\allowbreak flt.\allowbreak knn.\allowbreak avg.\allowbreak pcl} & carbon utilization; diauxic shift; fermentation; respiration & 20 & -- & 2 \\
\addlinespace[3pt]
200 & \textbf{Chitikila 2002}\newline \sourcetext{Chitikila C, Huisinga KL, Irvin JD, Basehoar AD, Pugh BF. Interplay of TBP inhibitors in global transcriptional control. \emph{Mol Cell}} 2002;10(4):871--82. & \href{https://doi.org/10.1016/s1097-2765(02)00683-4}{doi:10.\allowbreak 1016/\allowbreak s1097-\allowbreak 2765(02)00683-\allowbreak 4}\newline \href{https://pubmed.ncbi.nlm.nih.gov/12419230/}{PMID 12419230} & {\ttfamily\scriptsize 2010.\allowbreak Chitikila02.\allowbreak flt.\allowbreak knn.\allowbreak avg.\allowbreak pcl} & transcription & 20 & -- & 2 \\
\addlinespace[3pt]
209 & \textbf{Sabet 2004}\newline \sourcetext{Sabet N, Volo S, Yu C, Madigan JP, Morse RH. Genome-wide analysis of the relationship between transcriptional regulation by Rpd3p and the histone H3 and H4 amino termini in budding yeast. \emph{Mol Cell Biol}} 2004;24(20):8823--33. & \href{https://doi.org/10.1128/MCB.24.20.8823-8833.2004}{doi:10.\allowbreak 1128/\allowbreak MCB.\allowbreak 24.\allowbreak 20.\allowbreak 8823-\allowbreak 8833.\allowbreak 2004}\newline \href{https://pubmed.ncbi.nlm.nih.gov/15456858/}{PMID 15456858}\newline \href{https://pmc.ncbi.nlm.nih.gov/articles/PMC517902/}{PMC517902} & {\ttfamily\scriptsize 2010.\allowbreak Sabet04.\allowbreak filter.\allowbreak flt.\allowbreak knn.\allowbreak avg.\allowbreak div.\allowbreak log.\allowbreak pcl} & histone modification & 18 & 1 & 2 \\
\addlinespace[3pt]
212 & \textbf{Bernstein 2000}\newline \sourcetext{Bernstein BE, Tong JK, Schreiber SL. Genomewide studies of histone deacetylase function in yeast. \emph{Proc Natl Acad Sci U S A}} 2000;97(25):13708--13. & \href{https://doi.org/10.1073/pnas.250477697}{doi:10.\allowbreak 1073/\allowbreak pnas.\allowbreak 250477697}\newline \href{https://pubmed.ncbi.nlm.nih.gov/11095743/}{PMID 11095743}\newline \href{https://pmc.ncbi.nlm.nih.gov/articles/PMC17640/}{PMC17640} & {\ttfamily\scriptsize 2010.\allowbreak Bernstein00\_\allowbreak HDACrpd3sin3hda1.\allowbreak log.\allowbreak flt.\allowbreak knn.\allowbreak avg.\allowbreak pcl\newline 2010.\allowbreak Bernstein00\_\allowbreak HDACsin3sap30ume6hda1hos2hos3.\allowbreak filter.\allowbreak flt.\allowbreak knn.\allowbreak avg.\allowbreak div.\allowbreak log.\allowbreak pcl\newline 2010.\allowbreak Bernstein00\_\allowbreak TSA.\allowbreak filter.\allowbreak flt.\allowbreak knn.\allowbreak avg.\allowbreak div.\allowbreak log.\allowbreak pcl} & chromatin organization; histone modification & 18 & -- & 2 \\
\addlinespace[3pt]
233 & \textbf{Cho 1998}\newline \sourcetext{Cho RJ, Campbell MJ, Winzeler EA, et al., Davis RW. A genome-wide transcriptional analysis of the mitotic cell cycle. \emph{Mol Cell}} 1998;2(1):65--73. & \href{https://doi.org/10.1016/s1097-2765(00)80114-8}{doi:10.\allowbreak 1016/\allowbreak s1097-\allowbreak 2765(00)80114-\allowbreak 8}\newline \href{https://pubmed.ncbi.nlm.nih.gov/9702192/}{PMID 9702192} & {\ttfamily\scriptsize 2010.\allowbreak Cho98.\allowbreak filter.\allowbreak flt.\allowbreak knn.\allowbreak avg.\allowbreak div.\allowbreak log.\allowbreak pcl} & mitotic cell cycle & 17 & -- & 2 \\
\addlinespace[3pt]
234 & \textbf{Clark 2002}\newline \sourcetext{Clark TA, Sugnet CW, Ares M Jr. Genomewide analysis of mRNA processing in yeast using splicing-specific microarrays. \emph{Science}} 2002;296(5569):907--10. & \href{https://doi.org/10.1126/science.1069415}{doi:10.\allowbreak 1126/\allowbreak science.\allowbreak 1069415}\newline \href{https://pubmed.ncbi.nlm.nih.gov/11988574/}{PMID 11988574} & {\ttfamily\scriptsize 2010.\allowbreak Clark02\_\allowbreak orig.\allowbreak flt.\allowbreak knn.\allowbreak avg.\allowbreak pcl} & mRNA processing & 17 & -- & 2 \\
\addlinespace[3pt]
237 & \textbf{Haugen 2004}\newline \sourcetext{Haugen AC, Kelley R, Collins JB, et al., Van Houten B. Integrating phenotypic and expression profiles to map arsenic-response networks. \emph{Genome Biol}} 2004;5(12):R95. & \href{https://doi.org/10.1186/gb-2004-5-12-r95}{doi:10.\allowbreak 1186/\allowbreak gb-\allowbreak 2004-\allowbreak 5-\allowbreak 12-\allowbreak r95}\newline \href{https://pubmed.ncbi.nlm.nih.gov/15575969/}{PMID 15575969}\newline \href{https://pmc.ncbi.nlm.nih.gov/articles/PMC545798/}{PMC545798} & {\ttfamily\scriptsize 2010.\allowbreak arsenic\_\allowbreak time\_\allowbreak dose.\allowbreak pcl\newline 2010.\allowbreak arsenic\_\allowbreak deletions.\allowbreak pcl} & chemical stimulus; metal or metalloid ion stress & 17 & -- & 2 \\
\addlinespace[3pt]
240 & \textbf{Lopez 2000}\newline \sourcetext{L\'opez MC, Baker HV. Understanding the growth phenotype of the yeast gcr1 mutant in terms of global genomic expression patterns. \emph{J Bacteriol}} 2000;182(17):4970--8. & \href{https://doi.org/10.1128/JB.182.17.4970-4978.2000}{doi:10.\allowbreak 1128/\allowbreak JB.\allowbreak 182.\allowbreak 17.\allowbreak 4970-\allowbreak 4978.\allowbreak 2000}\newline \href{https://pubmed.ncbi.nlm.nih.gov/10940042/}{PMID 10940042}\newline \href{https://pmc.ncbi.nlm.nih.gov/articles/PMC111378/}{PMC111378} & {\ttfamily\scriptsize 2010.\allowbreak Lopez00.\allowbreak filter.\allowbreak flt.\allowbreak knn.\allowbreak avg.\allowbreak div.\allowbreak log.\allowbreak pcl} & carbon utilization; fermentation; respiration & 17 & -- & 2 \\
\addlinespace[3pt]
264 & \textbf{Carmel-Harel 2001}\newline \sourcetext{Carmel-Harel O, Stearman R, Gasch AP, Botstein D, Brown PO, Storz G. Role of thioredoxin reductase in the Yap1p-dependent response to oxidative stress in Saccharomyces cerevisiae. \emph{Mol Microbiol}} 2001;39(3):595--605. & \href{https://doi.org/10.1046/j.1365-2958.2001.02255.x}{doi:10.\allowbreak 1046/\allowbreak j.\allowbreak 1365-\allowbreak 2958.\allowbreak 2001.\allowbreak 02255.\allowbreak x}\newline \href{https://pubmed.ncbi.nlm.nih.gov/11169101/}{PMID 11169101} & {\ttfamily\scriptsize 2010.\allowbreak CarmelHarel01.\allowbreak flt.\allowbreak knn.\allowbreak avg.\allowbreak pcl} & chemical stimulus; oxidative stress; stress & 15 & -- & 2 \\
\addlinespace[3pt]
270 & \textbf{Roberts 2000}\newline \sourcetext{Roberts CJ, Nelson B, Marton MJ, et al., Friend SH. Signaling and circuitry of multiple MAPK pathways revealed by a matrix of global gene expression profiles. \emph{Science}} 2000;287(5454):873--80. & \href{https://doi.org/10.1126/science.287.5454.873}{doi:10.\allowbreak 1126/\allowbreak science.\allowbreak 287.\allowbreak 5454.\allowbreak 873}\newline \href{https://pubmed.ncbi.nlm.nih.gov/10657304/}{PMID 10657304} & {\ttfamily\scriptsize 2010.\allowbreak alpha\_\allowbreak time.\allowbreak pcl\newline 2010.\allowbreak alpha\_\allowbreak conc.\allowbreak pcl} & cell morphogenesis; mating; signaling & 15 & -- & 2 \\
\addlinespace[3pt]
281 & \textbf{Hardwick 1999}\newline \sourcetext{Hardwick JS, Kuruvilla FG, Tong JK, Shamji AF, Schreiber SL. Rapamycin-modulated transcription defines the subset of nutrient-sensitive signaling pathways directly controlled by the Tor proteins. \emph{Proc Natl Acad Sci U S A}} 1999;96(26):14866--70. & \href{https://doi.org/10.1073/pnas.96.26.14866}{doi:10.\allowbreak 1073/\allowbreak pnas.\allowbreak 96.\allowbreak 26.\allowbreak 14866}\newline \href{https://pubmed.ncbi.nlm.nih.gov/10611304/}{PMID 10611304}\newline \href{https://pmc.ncbi.nlm.nih.gov/articles/PMC24739/}{PMC24739} & {\ttfamily\scriptsize 2010.\allowbreak Hardwick00.\allowbreak flt.\allowbreak knn.\allowbreak avg.\allowbreak pcl} & amino acid utilization; chemical stimulus; signaling & 14 & -- & 2 \\
\addlinespace[3pt]
293 & \textbf{Leber 2004}\newline \sourcetext{Leber JH, Bernales S, Walter P. IRE1-independent gain control of the unfolded protein response. \emph{PLoS Biol}} 2004;2(8):E235. & \href{https://doi.org/10.1371/journal.pbio.0020235}{doi:10.\allowbreak 1371/\allowbreak journal.\allowbreak pbio.\allowbreak 0020235}\newline \href{https://pubmed.ncbi.nlm.nih.gov/15314654/}{PMID 15314654}\newline \href{https://pmc.ncbi.nlm.nih.gov/articles/PMC509300/}{PMC509300} & {\ttfamily\scriptsize 2010.\allowbreak Leber04.\allowbreak flt.\allowbreak knn.\allowbreak avg.\allowbreak pcl} & chemical stimulus; response to unfolded protein; stress & 13 & -- & 1 \\
\addlinespace[3pt]
297 & \textbf{Carroll 2001}\newline \sourcetext{Carroll AS, Bishop AC, DeRisi JL, Shokat KM, O'Shea EK. Chemical inhibition of the Pho85 cyclin-dependent kinase reveals a role in the environmental stress response. \emph{Proc Natl Acad Sci U S A}} 2001;98(22):12578--83. & \href{https://doi.org/10.1073/pnas.211195798}{doi:10.\allowbreak 1073/\allowbreak pnas.\allowbreak 211195798}\newline \href{https://pubmed.ncbi.nlm.nih.gov/11675494/}{PMID 11675494}\newline \href{https://pmc.ncbi.nlm.nih.gov/articles/PMC60096/}{PMC60096} & {\ttfamily\scriptsize 2010.\allowbreak Carroll01.\allowbreak log.\allowbreak flt.\allowbreak knn.\allowbreak avg.\allowbreak pcl} & protein phosphorylation; stress & 12 & 2 & 2 \\
\addlinespace[3pt]
299 & \textbf{Boer 2005}\newline \sourcetext{Boer VM, Daran JM, Almering MJ, de Winde JH, Pronk JT. Contribution of the Saccharomyces cerevisiae transcriptional regulator Leu3p to physiology and gene expression in nitrogen- and carbon-limited chemostat cultures. \emph{FEMS Yeast Res}} 2005;5(10):885--97. & \href{https://doi.org/10.1016/j.femsyr.2005.04.003}{doi:10.\allowbreak 1016/\allowbreak j.\allowbreak femsyr.\allowbreak 2005.\allowbreak 04.\allowbreak 003}\newline \href{https://pubmed.ncbi.nlm.nih.gov/15949974/}{PMID 15949974} & {\ttfamily\scriptsize 2010.\allowbreak Boer05.\allowbreak filter.\allowbreak flt.\allowbreak knn.\allowbreak avg.\allowbreak div.\allowbreak log.\allowbreak pcl} & carbon utilization; nitrogen utilization; transcription & 12 & 1 & 2 \\
\addlinespace[3pt]
321 & \textbf{De Nadal 2004}\newline \sourcetext{De Nadal E, Zapater M, Alepuz PM, Sumoy L, Mas G, Posas F. The MAPK Hog1 recruits Rpd3 histone deacetylase to activate osmoresponsive genes. \emph{Nature}} 2004;427(6972):370--4. & \href{https://doi.org/10.1038/nature02258}{doi:10.\allowbreak 1038/\allowbreak nature02258}\newline \href{https://pubmed.ncbi.nlm.nih.gov/14737171/}{PMID 14737171} & {\ttfamily\scriptsize 2010.\allowbreak deNadal04.\allowbreak flt.\allowbreak knn.\allowbreak avg.\allowbreak pcl} & histone modification; osmotic stress; stress & 12 & -- & 2 \\
\addlinespace[3pt]
334 & \textbf{Martin 2004}\newline \sourcetext{Martin DE, Demougin P, Hall MN, Bellis M. Rank Difference Analysis of Microarrays (RDAM), a novel approach to statistical analysis of microarray expression profiling data. \emph{BMC Bioinformatics}} 2004;5:148. & \href{https://doi.org/10.1186/1471-2105-5-148}{doi:10.\allowbreak 1186/\allowbreak 1471-\allowbreak 2105-\allowbreak 5-\allowbreak 148}\newline \href{https://pubmed.ncbi.nlm.nih.gov/15476558/}{PMID 15476558}\newline \href{https://pmc.ncbi.nlm.nih.gov/articles/PMC526220/}{PMC526220} & {\ttfamily\scriptsize 2010.\allowbreak Martin04.\allowbreak filter.\allowbreak flt.\allowbreak knn.\allowbreak avg.\allowbreak div.\allowbreak log.\allowbreak pcl} & protein phosphorylation; signaling & 12 & -- & 2 \\
\addlinespace[3pt]
337 & \textbf{Olesen 2002}\newline \sourcetext{Olesen K, Felding T, Gjermansen C, Hansen J. The dynamics of the Saccharomyces carlsbergensis brewing yeast transcriptome during a production-scale lager beer fermentation. \emph{FEMS Yeast Res}} 2002;2(4):563--73. & \href{https://doi.org/10.1111/j.1567-1364.2002.tb00123.x}{doi:10.\allowbreak 1111/\allowbreak j.\allowbreak 1567-\allowbreak 1364.\allowbreak 2002.\allowbreak tb00123.\allowbreak x}\newline \href{https://pubmed.ncbi.nlm.nih.gov/12702272/}{PMID 12702272} & {\ttfamily\scriptsize 2010.\allowbreak Olesen02.\allowbreak filter.\allowbreak flt.\allowbreak knn.\allowbreak avg.\allowbreak div.\allowbreak log.\allowbreak pcl} & carbon utilization; fermentation & 12 & -- & 2 \\
\addlinespace[3pt]
342 & \textbf{Robertson 2000}\newline \sourcetext{Robertson LS, Causton HC, Young RA, Fink GR. The yeast A kinases differentially regulate iron uptake and respiratory function. \emph{Proc Natl Acad Sci U S A}} 2000;97(11):5984--8. & \href{https://doi.org/10.1073/pnas.100113397}{doi:10.\allowbreak 1073/\allowbreak pnas.\allowbreak 100113397}\newline \href{https://pubmed.ncbi.nlm.nih.gov/10811893/}{PMID 10811893}\newline \href{https://pmc.ncbi.nlm.nih.gov/articles/PMC18545/}{PMC18545} & {\ttfamily\scriptsize 2010.\allowbreak Robertson00.\allowbreak filter.\allowbreak flt.\allowbreak knn.\allowbreak avg.\allowbreak div.\allowbreak log.\allowbreak pcl} & protein phosphorylation; signaling & 12 & -- & 2 \\
\addlinespace[3pt]
347 & \textbf{Sudarsanam 2000}\newline \sourcetext{Sudarsanam P, Iyer VR, Brown PO, Winston F. Whole-genome expression analysis of snf/swi mutants of Saccharomyces cerevisiae. \emph{Proc Natl Acad Sci U S A}} 2000;97(7):3364--9. & \href{https://doi.org/10.1073/pnas.97.7.3364}{doi:10.\allowbreak 1073/\allowbreak pnas.\allowbreak 97.\allowbreak 7.\allowbreak 3364}\newline \href{https://pubmed.ncbi.nlm.nih.gov/10725359/}{PMID 10725359}\newline \href{https://pmc.ncbi.nlm.nih.gov/articles/PMC16245/}{PMC16245} & {\ttfamily\scriptsize 2010.\allowbreak snf\_\allowbreak swi\_\allowbreak mutants.\allowbreak pcl} & chromatin organization & 12 & -- & 2 \\
\addlinespace[3pt]
350 & \textbf{Takagi 2005}\newline \sourcetext{Takagi Y, Masuda CA, Chang WH, et al., Kornberg RD. Ubiquitin ligase activity of TFIIH and the transcriptional response to DNA damage. \emph{Mol Cell}} 2005;18(2):237--43. & \href{https://doi.org/10.1016/j.molcel.2005.03.007}{doi:10.\allowbreak 1016/\allowbreak j.\allowbreak molcel.\allowbreak 2005.\allowbreak 03.\allowbreak 007}\newline \href{https://pubmed.ncbi.nlm.nih.gov/15837426/}{PMID 15837426} & {\ttfamily\scriptsize 2010.\allowbreak Takagi05.\allowbreak filter.\allowbreak flt.\allowbreak knn.\allowbreak avg.\allowbreak div.\allowbreak log.\allowbreak pcl} & DNA damage stimulus; transcription; ubiquitin or ULP modification & 12 & -- & 2 \\
\addlinespace[3pt]
360 & \textbf{Bulik 2003}\newline \sourcetext{Bulik DA, Olczak M, Lucero HA, Osmond BC, Robbins PW, Specht CA. Chitin synthesis in Saccharomyces cerevisiae in response to supplementation of growth medium with glucosamine and cell wall stress. \emph{Eukaryot Cell}} 2003;2(5):886--900. & \href{https://doi.org/10.1128/EC.2.5.886-900.2003}{doi:10.\allowbreak 1128/\allowbreak EC.\allowbreak 2.\allowbreak 5.\allowbreak 886-\allowbreak 900.\allowbreak 2003}\newline \href{https://pubmed.ncbi.nlm.nih.gov/14555471/}{PMID 14555471}\newline \href{https://pmc.ncbi.nlm.nih.gov/articles/PMC219353/}{PMC219353} & {\ttfamily\scriptsize 2010.\allowbreak Bulik03.\allowbreak filter.\allowbreak flt.\allowbreak knn.\allowbreak avg.\allowbreak div.\allowbreak log.\allowbreak pcl} & cell wall organization; chemical stimulus; stress & 11 & -- & 2 \\
\addlinespace[3pt]
361 & \textbf{Cohen 2002}\newline \sourcetext{Cohen BA, Pilpel Y, Mitra RD, Church GM. Discrimination between paralogs using microarray analysis: application to the Yap1p and Yap2p transcriptional networks. \emph{Mol Biol Cell}} 2002;13(5):1608--14. & \href{https://doi.org/10.1091/mbc.01-10-0472}{doi:10.\allowbreak 1091/\allowbreak mbc.\allowbreak 01-\allowbreak 10-\allowbreak 0472}\newline \href{https://pubmed.ncbi.nlm.nih.gov/12006656/}{PMID 12006656}\newline \href{https://pmc.ncbi.nlm.nih.gov/articles/PMC111130/}{PMC111130} & {\ttfamily\scriptsize 2010.\allowbreak Cohen02.\allowbreak filter.\allowbreak flt.\allowbreak knn.\allowbreak avg.\allowbreak div.\allowbreak log.\allowbreak pcl} & chemical stimulus; oxidative stress; stress; transcription & 11 & -- & 2 \\
\addlinespace[3pt]
362 & \textbf{Epstein 2001}\newline \sourcetext{Epstein CB, Waddle JA, Hale W 4th, et al., Butow RA. Genome-wide responses to mitochondrial dysfunction. \emph{Mol Biol Cell}} 2001;12(2):297--308. & \href{https://doi.org/10.1091/mbc.12.2.297}{doi:10.\allowbreak 1091/\allowbreak mbc.\allowbreak 12.\allowbreak 2.\allowbreak 297}\newline \href{https://pubmed.ncbi.nlm.nih.gov/11179416/}{PMID 11179416}\newline \href{https://pmc.ncbi.nlm.nih.gov/articles/PMC30944/}{PMC30944} & {\ttfamily\scriptsize 2010.\allowbreak Epstein00.\allowbreak flt.\allowbreak knn.\allowbreak avg.\allowbreak pcl} & carbon utilization; respiration & 11 & -- & 2 \\
\addlinespace[3pt]
364 & \textbf{Galitski 1999}\newline \sourcetext{Galitski T, Saldanha AJ, Styles CA, Lander ES, Fink GR. Ploidy regulation of gene expression. \emph{Science}} 1999;285(5425):251--4. & \href{https://doi.org/10.1126/science.285.5425.251}{doi:10.\allowbreak 1126/\allowbreak science.\allowbreak 285.\allowbreak 5425.\allowbreak 251}\newline \href{https://pubmed.ncbi.nlm.nih.gov/10398601/}{PMID 10398601} & {\ttfamily\scriptsize 2010.\allowbreak ploidy.\allowbreak pcl} & cell morphogenesis; mitotic cell cycle & 11 & -- & 2 \\
\addlinespace[3pt]
366 & \textbf{Madhani 1999}\newline \sourcetext{Madhani HD, Galitski T, Lander ES, Fink GR. Effectors of a developmental mitogen-activated protein kinase cascade revealed by expression signatures of signaling mutants. \emph{Proc Natl Acad Sci U S A}} 1999;96(22):12530--5. & \href{https://doi.org/10.1073/pnas.96.22.12530}{doi:10.\allowbreak 1073/\allowbreak pnas.\allowbreak 96.\allowbreak 22.\allowbreak 12530}\newline \href{https://pubmed.ncbi.nlm.nih.gov/10535956/}{PMID 10535956}\newline \href{https://pmc.ncbi.nlm.nih.gov/articles/PMC22972/}{PMC22972} & {\ttfamily\scriptsize 2010.\allowbreak Madhani99.\allowbreak filter.\allowbreak flt.\allowbreak knn.\allowbreak avg.\allowbreak div.\allowbreak log.\allowbreak pcl} & filamentous growth; signaling & 11 & -- & 2 \\
\addlinespace[3pt]
378 & \textbf{Huang 2004}\newline \sourcetext{Huang J, Zhu H, Haggarty SJ, et al., Schreiber SL. Finding new components of the target of rapamycin (TOR) signaling network through chemical genetics and proteome chips. \emph{Proc Natl Acad Sci U S A}} 2004;101(47):16594--9. & \href{https://doi.org/10.1073/pnas.0407117101}{doi:10.\allowbreak 1073/\allowbreak pnas.\allowbreak 0407117101}\newline \href{https://pubmed.ncbi.nlm.nih.gov/15539461/}{PMID 15539461}\newline \href{https://pmc.ncbi.nlm.nih.gov/articles/PMC527135/}{PMC527135} & {\ttfamily\scriptsize 2010.\allowbreak rapamycin.\allowbreak pcl} & chemical stimulus; signaling & 10 & -- & 2 \\
\addlinespace[3pt]
382 & \textbf{Prinz 2004}\newline \sourcetext{Prinz S, Avila-Campillo I, Aldridge C, et al., Galitski T. Control of yeast filamentous-form growth by modules in an integrated molecular network. \emph{Genome Res}} 2004;14(3):380--90. & \href{https://doi.org/10.1101/gr.2020604}{doi:10.\allowbreak 1101/\allowbreak gr.\allowbreak 2020604}\newline \href{https://pubmed.ncbi.nlm.nih.gov/14993204/}{PMID 14993204}\newline \href{https://pmc.ncbi.nlm.nih.gov/articles/PMC353223/}{PMC353223} & {\ttfamily\scriptsize 2010.\allowbreak Prinz04.\allowbreak flt.\allowbreak knn.\allowbreak avg.\allowbreak pcl} & filamentous growth & 10 & -- & 2 \\
\addlinespace[3pt]
384 & \textbf{Travers 2000}\newline \sourcetext{Travers KJ, Patil CK, Wodicka L, Lockhart DJ, Weissman JS, Walter P. Functional and genomic analyses reveal an essential coordination between the unfolded protein response and ER-associated degradation. \emph{Cell}} 2000;101(3):249--58. & \href{https://doi.org/10.1016/s0092-8674(00)80835-1}{doi:10.\allowbreak 1016/\allowbreak s0092-\allowbreak 8674(00)80835-\allowbreak 1}\newline \href{https://pubmed.ncbi.nlm.nih.gov/10847680/}{PMID 10847680} & {\ttfamily\scriptsize 2010.\allowbreak unfolded\_\allowbreak protein\_\allowbreak response.\allowbreak pcl} & chemical stimulus; response to unfolded protein; stress & 10 & -- & 2 \\
\addlinespace[3pt]
399 & \textbf{Lyons 2000}\newline \sourcetext{Lyons TJ, Gasch AP, Gaither LA, Botstein D, Brown PO, Eide DJ. Genome-wide characterization of the Zap1p zinc-responsive regulon in yeast. \emph{Proc Natl Acad Sci U S A}} 2000;97(14):7957--62. & \href{https://doi.org/10.1073/pnas.97.14.7957}{doi:10.\allowbreak 1073/\allowbreak pnas.\allowbreak 97.\allowbreak 14.\allowbreak 7957}\newline \href{https://pubmed.ncbi.nlm.nih.gov/10884426/}{PMID 10884426}\newline \href{https://pmc.ncbi.nlm.nih.gov/articles/PMC16652/}{PMC16652} & {\ttfamily\scriptsize 2010.\allowbreak Lyons00.\allowbreak flt.\allowbreak knn.\allowbreak avg.\allowbreak pcl} & cellular ion homeostasis; metal or metalloid ion stress & 9 & -- & 2 \\
\addlinespace[3pt]
408 & \textbf{Angus-Hill 2001}\newline \sourcetext{Angus-Hill ML, Schlichter A, Roberts D, Erdjument-Bromage H, Tempst P, Cairns BR. A Rsc3/Rsc30 zinc cluster dimer reveals novel roles for the chromatin remodeler RSC in gene expression and cell cycle control. \emph{Mol Cell}} 2001;7(4):741--51. & \href{https://doi.org/10.1016/s1097-2765(01)00219-2}{doi:10.\allowbreak 1016/\allowbreak s1097-\allowbreak 2765(01)00219-\allowbreak 2}\newline \href{https://pubmed.ncbi.nlm.nih.gov/11336698/}{PMID 11336698} & {\ttfamily\scriptsize 2010.\allowbreak Angus01.\allowbreak flt.\allowbreak knn.\allowbreak avg.\allowbreak pcl} & chromatin organization & 8 & 1 & 2 \\
\addlinespace[3pt]
421 & \textbf{Fry 2003}\newline \sourcetext{Fry RC, Sambandan TG, Rha C. DNA damage and stress transcripts in Saccharomyces cerevisiae mutant sgs1. \emph{Mech Ageing Dev}} 2003;124(7):839--46. & \href{https://doi.org/10.1016/s0047-6374(03)00144-1}{doi:10.\allowbreak 1016/\allowbreak s0047-\allowbreak 6374(03)00144-\allowbreak 1}\newline \href{https://pubmed.ncbi.nlm.nih.gov/12875747/}{PMID 12875747} & {\ttfamily\scriptsize 2010.\allowbreak Fry03.\allowbreak filter.\allowbreak flt.\allowbreak knn.\allowbreak avg.\allowbreak div.\allowbreak log.\allowbreak pcl} & DNA damage stimulus; stress & 8 & -- & 2 \\
\addlinespace[3pt]
428 & \textbf{Lee 2000}\newline \sourcetext{Lee SE, Pellicioli A, Demeter J, et al., Haber JE. Arrest, adaptation, and recovery following a chromosome double-strand break in Saccharomyces cerevisiae. \emph{Cold Spring Harb Symp Quant Biol}} 2000;65:303--14. & \href{https://doi.org/10.1101/sqb.2000.65.303}{doi:10.\allowbreak 1101/\allowbreak sqb.\allowbreak 2000.\allowbreak 65.\allowbreak 303}\newline \href{https://pubmed.ncbi.nlm.nih.gov/12760044/}{PMID 12760044} & {\ttfamily\scriptsize 2010.\allowbreak HO\_\allowbreak endonuclease.\allowbreak pcl} & DNA damage stimulus & 8 & -- & 2 \\
\addlinespace[3pt]
433 & \textbf{Ogawa 2000}\newline \sourcetext{Ogawa N, DeRisi J, Brown PO. New components of a system for phosphate accumulation and polyphosphate metabolism in Saccharomyces cerevisiae revealed by genomic expression analysis. \emph{Mol Biol Cell}} 2000;11(12):4309--21. & \href{https://doi.org/10.1091/mbc.11.12.4309}{doi:10.\allowbreak 1091/\allowbreak mbc.\allowbreak 11.\allowbreak 12.\allowbreak 4309}\newline \href{https://pubmed.ncbi.nlm.nih.gov/11102525/}{PMID 11102525}\newline \href{https://pmc.ncbi.nlm.nih.gov/articles/PMC15074/}{PMC15074} & {\ttfamily\scriptsize 2010.\allowbreak Ogawa00\_\allowbreak set1.\allowbreak flt.\allowbreak knn.\allowbreak avg.\allowbreak pcl\newline 2010.\allowbreak Ogawa00\_\allowbreak set2.\allowbreak flt.\allowbreak knn.\allowbreak avg.\allowbreak pcl} & phosphorus utilization; signaling; starvation & 8 & -- & 2 \\
\addlinespace[3pt]
443 & \textbf{Smith 2002}\newline \sourcetext{Smith JJ, Marelli M, Christmas RH, et al., Aitchison JD. Transcriptome profiling to identify genes involved in peroxisome assembly and function. \emph{J Cell Biol}} 2002;158(2):259--71. & \href{https://doi.org/10.1083/jcb.200204059}{doi:10.\allowbreak 1083/\allowbreak jcb.\allowbreak 200204059}\newline \href{https://pubmed.ncbi.nlm.nih.gov/12135984/}{PMID 12135984}\newline \href{https://pmc.ncbi.nlm.nih.gov/articles/PMC2173120/}{PMC2173120} & {\ttfamily\scriptsize 2010.\allowbreak peroxisome.\allowbreak pcl} & carbon utilization; lipid metabolism & 8 & -- & 2 \\
\addlinespace[3pt]
447 & \textbf{Williams 2002}\newline \sourcetext{Williams RM, Primig M, Washburn BK, et al., Esposito RE. The Ume6 regulon coordinates metabolic and meiotic gene expression in yeast. \emph{Proc Natl Acad Sci U S A}} 2002;99(21):13431--6. & \href{https://doi.org/10.1073/pnas.202495299}{doi:10.\allowbreak 1073/\allowbreak pnas.\allowbreak 202495299}\newline \href{https://pubmed.ncbi.nlm.nih.gov/12370439/}{PMID 12370439}\newline \href{https://pmc.ncbi.nlm.nih.gov/articles/PMC129690/}{PMC129690} & {\ttfamily\scriptsize 2010.\allowbreak Williams02.\allowbreak filter.\allowbreak flt.\allowbreak knn.\allowbreak avg.\allowbreak div.\allowbreak log.\allowbreak pcl} & carbon utilization; sporulation & 8 & -- & 2 \\
\addlinespace[3pt]
452 & \textbf{Bedalov 2001}\newline \sourcetext{Bedalov A, Gatbonton T, Irvine WP, Gottschling DE, Simon JA. Identification of a small molecule inhibitor of Sir2p. \emph{Proc Natl Acad Sci U S A}} 2001;98(26):15113--8. & \href{https://doi.org/10.1073/pnas.261574398}{doi:10.\allowbreak 1073/\allowbreak pnas.\allowbreak 261574398}\newline \href{https://pubmed.ncbi.nlm.nih.gov/11752457/}{PMID 11752457}\newline \href{https://pmc.ncbi.nlm.nih.gov/articles/PMC64992/}{PMC64992} & {\ttfamily\scriptsize 2010.\allowbreak Bedalov01.\allowbreak flt.\allowbreak knn.\allowbreak avg.\allowbreak pcl} & chemical stimulus; chromatin organization; histone modification & 7 & -- & 2 \\
\addlinespace[3pt]
453 & \textbf{Bro 2003}\newline \sourcetext{Bro C, Regenberg B, Lagniel G, Labarre J, Montero-Lomel\'{\i} M, Nielsen J. Transcriptional, proteomic, and metabolic responses to lithium in galactose-grown yeast cells. \emph{J Biol Chem}} 2003;278(34):32141--9. & \href{https://doi.org/10.1074/jbc.M304478200}{doi:10.\allowbreak 1074/\allowbreak jbc.\allowbreak M304478200}\newline \href{https://pubmed.ncbi.nlm.nih.gov/12791685/}{PMID 12791685} & {\ttfamily\scriptsize 2010.\allowbreak Bro03.\allowbreak filter.\allowbreak flt.\allowbreak knn.\allowbreak avg.\allowbreak div.\allowbreak log.\allowbreak pcl} & chemical stimulus; metal or metalloid ion stress & 7 & -- & 2 \\
\addlinespace[3pt]
455 & \textbf{Chu 1998}\newline \sourcetext{Chu S, DeRisi J, Eisen M, et al., Herskowitz I. The transcriptional program of sporulation in budding yeast. \emph{Science}} 1998;282(5389):699--705. & \href{https://doi.org/10.1126/science.282.5389.699}{doi:10.\allowbreak 1126/\allowbreak science.\allowbreak 282.\allowbreak 5389.\allowbreak 699}\newline \href{https://pubmed.ncbi.nlm.nih.gov/9784122/}{PMID 9784122} & {\ttfamily\scriptsize 2010.\allowbreak Chu98.\allowbreak flt.\allowbreak knn.\allowbreak avg.\allowbreak pcl} & sporulation & 7 & -- & 2 \\
\addlinespace[3pt]
456 & \textbf{DeRisi 1997}\newline \sourcetext{DeRisi JL, Iyer VR, Brown PO. Exploring the metabolic and genetic control of gene expression on a genomic scale. \emph{Science}} 1997;278(5338):680--6. & \href{https://doi.org/10.1126/science.278.5338.680}{doi:10.\allowbreak 1126/\allowbreak science.\allowbreak 278.\allowbreak 5338.\allowbreak 680}\newline \href{https://pubmed.ncbi.nlm.nih.gov/9381177/}{PMID 9381177} & {\ttfamily\scriptsize 2010.\allowbreak diauxic.\allowbreak pcl} & carbon utilization; diauxic shift; fermentation; respiration & 7 & -- & 2 \\
\addlinespace[3pt]
460 & \textbf{Marton 1998}\newline \sourcetext{Marton MJ, DeRisi JL, Bennett HA, et al., Friend SH. Drug target validation and identification of secondary drug target effects using DNA microarrays. \emph{Nat Med}} 1998;4(11):1293--301. & \href{https://doi.org/10.1038/3282}{doi:10.\allowbreak 1038/\allowbreak 3282}\newline \href{https://pubmed.ncbi.nlm.nih.gov/9809554/}{PMID 9809554} & {\ttfamily\scriptsize 2010.\allowbreak Marton98.\allowbreak log.\allowbreak flt.\allowbreak knn.\allowbreak avg.\allowbreak pcl} & chemical stimulus; protein dephosphorylation; signaling & 7 & -- & 2 \\
\addlinespace[3pt]
461 & \textbf{Wyrick 1999}\newline \sourcetext{Wyrick JJ, Holstege FC, Jennings EG, et al., Young RA. Chromosomal landscape of nucleosome-dependent gene expression and silencing in yeast. \emph{Nature}} 1999;402(6760):418--21. & \href{https://doi.org/10.1038/46567}{doi:10.\allowbreak 1038/\allowbreak 46567}\newline \href{https://pubmed.ncbi.nlm.nih.gov/10586882/}{PMID 10586882} & {\ttfamily\scriptsize 2010.\allowbreak histone.\allowbreak pcl} & chromatin organization & 7 & -- & 2 \\
\addlinespace[3pt]
462 & \textbf{Yeang 2005}\newline \sourcetext{Yeang CH, Mak HC, McCuine S, Workman C, Jaakkola T, Ideker T. Validation and refinement of gene-regulatory pathways on a network of physical interactions. \emph{Genome Biol}} 2005;6(7):R62. & \href{https://doi.org/10.1186/gb-2005-6-7-r62}{doi:10.\allowbreak 1186/\allowbreak gb-\allowbreak 2005-\allowbreak 6-\allowbreak 7-\allowbreak r62}\newline \href{https://pubmed.ncbi.nlm.nih.gov/15998451/}{PMID 15998451}\newline \href{https://pmc.ncbi.nlm.nih.gov/articles/PMC1175993/}{PMC1175993} & {\ttfamily\scriptsize 2010.\allowbreak Yeang05.\allowbreak flt.\allowbreak knn.\allowbreak avg.\allowbreak pcl} & transcription & 7 & -- & 2 \\
\addlinespace[3pt]
466 & \textbf{Rodriguez-Navarro 2004}\newline \sourcetext{Rodr\'{\i}guez-Navarro S, Fischer T, Luo MJ, et al., Hurt E. Sus1, a functional component of the SAGA histone acetylase complex and the nuclear pore-associated mRNA export machinery. \emph{Cell}} 2004;116(1):75--86. & \href{https://doi.org/10.1016/s0092-8674(03)01025-0}{doi:10.\allowbreak 1016/\allowbreak s0092-\allowbreak 8674(03)01025-\allowbreak 0}\newline \href{https://pubmed.ncbi.nlm.nih.gov/14718168/}{PMID 14718168} & {\ttfamily\scriptsize 2010.\allowbreak RodriguezNavarro04.\allowbreak filter.\allowbreak flt.\allowbreak knn.\allowbreak avg.\allowbreak div.\allowbreak log.\allowbreak pcl} & histone modification & 6 & 1 & 2 \\
\addlinespace[3pt]
479 & \textbf{Gross 2000}\newline \sourcetext{Gross C, Kelleher M, Iyer VR, Brown PO, Winge DR. Identification of the copper regulon in Saccharomyces cerevisiae by DNA microarrays. \emph{J Biol Chem}} 2000;275(41):32310--6. & \href{https://doi.org/10.1074/jbc.M005946200}{doi:10.\allowbreak 1074/\allowbreak jbc.\allowbreak M005946200}\newline \href{https://pubmed.ncbi.nlm.nih.gov/10922376/}{PMID 10922376} & {\ttfamily\scriptsize 2010.\allowbreak Gross00.\allowbreak flt.\allowbreak knn.\allowbreak avg.\allowbreak pcl} & cellular ion homeostasis; metal or metalloid ion stress & 6 & -- & 2 \\
\addlinespace[3pt]
487 & \textbf{Jin 2004}\newline \sourcetext{Jin YS, Laplaza JM, Jeffries TW. Saccharomyces cerevisiae engineered for xylose metabolism exhibits a respiratory response. \emph{Appl Environ Microbiol}} 2004;70(11):6816--25. & \href{https://doi.org/10.1128/AEM.70.11.6816-6825.2004}{doi:10.\allowbreak 1128/\allowbreak AEM.\allowbreak 70.\allowbreak 11.\allowbreak 6816-\allowbreak 6825.\allowbreak 2004}\newline \href{https://pubmed.ncbi.nlm.nih.gov/15528549/}{PMID 15528549}\newline \href{https://pmc.ncbi.nlm.nih.gov/articles/PMC525251/}{PMC525251} & {\ttfamily\scriptsize 2010.\allowbreak Jin04.\allowbreak filter.\allowbreak flt.\allowbreak knn.\allowbreak avg.\allowbreak div.\allowbreak log.\allowbreak pcl} & carbon utilization; oxygen level alteration; respiration & 6 & -- & 2 \\
\addlinespace[3pt]
510 & \textbf{Shakoury-Elizeh 2004}\newline \sourcetext{Shakoury-Elizeh M, Tiedeman J, Rashford J, et al., Philpott CC. Transcriptional remodeling in response to iron deprivation in Saccharomyces cerevisiae. \emph{Mol Biol Cell}} 2004;15(3):1233--43. & \href{https://doi.org/10.1091/mbc.e03-09-0642}{doi:10.\allowbreak 1091/\allowbreak mbc.\allowbreak e03-\allowbreak 09-\allowbreak 0642}\newline \href{https://pubmed.ncbi.nlm.nih.gov/14668481/}{PMID 14668481}\newline \href{https://pmc.ncbi.nlm.nih.gov/articles/PMC363115/}{PMC363115} & {\ttfamily\scriptsize 2010.\allowbreak ShakouryElizeh04.\allowbreak flt.\allowbreak knn.\allowbreak avg.\allowbreak pcl} & metal or metalloid ion stress; starvation & 6 & -- & 2 \\
\addlinespace[3pt]
525 & \textbf{Baetz 2001}\newline \sourcetext{Baetz K, Moffat J, Haynes J, Chang M, Andrews B. Transcriptional coregulation by the cell integrity mitogen-activated protein kinase Slt2 and the cell cycle regulator Swi4. \emph{Mol Cell Biol}} 2001;21(19):6515--28. & \href{https://doi.org/10.1128/MCB.21.19.6515-6528.2001}{doi:10.\allowbreak 1128/\allowbreak MCB.\allowbreak 21.\allowbreak 19.\allowbreak 6515-\allowbreak 6528.\allowbreak 2001}\newline \href{https://pubmed.ncbi.nlm.nih.gov/11533240/}{PMID 11533240}\newline \href{https://pmc.ncbi.nlm.nih.gov/articles/PMC99798/}{PMC99798} & {\ttfamily\scriptsize 2010.\allowbreak Baetz01.\allowbreak flt.\allowbreak knn.\allowbreak avg.\allowbreak pcl} & mitotic cell cycle; signaling & 5 & 4 & 2 \\
\addlinespace[3pt]
529 & \textbf{Jones 2003}\newline \sourcetext{Jones DL, Petty J, Hoyle DC, et al., Stateva LI. Transcriptome profiling of a Saccharomyces cerevisiae mutant with a constitutively activated Ras/cAMP pathway. \emph{Physiol Genomics}} 2003;16(1):107--18. & \href{https://doi.org/10.1152/physiolgenomics.00139.2003}{doi:10.\allowbreak 1152/\allowbreak physiolgenomics.\allowbreak 00139.\allowbreak 2003}\newline \href{https://pubmed.ncbi.nlm.nih.gov/14570984/}{PMID 14570984} & {\ttfamily\scriptsize 2010.\allowbreak Jones03\_\allowbreak orig.\allowbreak flt.\allowbreak knn.\allowbreak avg.\allowbreak pcl} & signaling & 5 & -- & 2 \\
\addlinespace[3pt]
531 & \textbf{Roberts 2006}\newline \sourcetext{Roberts GG, Hudson AP. Transcriptome profiling of Saccharomyces cerevisiae during a transition from fermentative to glycerol-based respiratory growth reveals extensive metabolic and structural remodeling. \emph{Mol Genet Genomics}} 2006;276(2):170--86. & \href{https://doi.org/10.1007/s00438-006-0133-9}{doi:10.\allowbreak 1007/\allowbreak s00438-\allowbreak 006-\allowbreak 0133-\allowbreak 9}\newline \href{https://pubmed.ncbi.nlm.nih.gov/16741729/}{PMID 16741729} & {\ttfamily\scriptsize 2010.\allowbreak glucose\_\allowbreak glycerol\_\allowbreak transition.\allowbreak pcl} & carbon utilization; fermentation; respiration & 5 & -- & 2 \\
\addlinespace[3pt]
533 & \textbf{van Attikum 2004}\newline \sourcetext{van Attikum H, Fritsch O, Hohn B, Gasser SM. Recruitment of the INO80 complex by H2A phosphorylation links ATP-dependent chromatin remodeling with DNA double-strand break repair. \emph{Cell}} 2004;119(6):777--88. & \href{https://doi.org/10.1016/j.cell.2004.11.033}{doi:10.\allowbreak 1016/\allowbreak j.\allowbreak cell.\allowbreak 2004.\allowbreak 11.\allowbreak 033}\newline \href{https://pubmed.ncbi.nlm.nih.gov/15607975/}{PMID 15607975} & {\ttfamily\scriptsize 2010.\allowbreak Ino80Arp8.\allowbreak pcl\newline 2010.\allowbreak MMSresponse.\allowbreak pcl} & DNA damage stimulus; chromatin organization & 5 & -- & 2 \\
\addlinespace[3pt]
534 & \textbf{Orlandi 2004}\newline \sourcetext{Orlandi I, Bettiga M, Alberghina L, Vai M. Transcriptional profiling of ubp10 null mutant reveals altered subtelomeric gene expression and insurgence of oxidative stress response. \emph{J Biol Chem}} 2004;279(8):6414--25. & \href{https://doi.org/10.1074/jbc.M306464200}{doi:10.\allowbreak 1074/\allowbreak jbc.\allowbreak M306464200}\newline \href{https://pubmed.ncbi.nlm.nih.gov/14623890/}{PMID 14623890} & {\ttfamily\scriptsize 2010.\allowbreak Orlandi04.\allowbreak filter.\allowbreak flt.\allowbreak knn.\allowbreak avg.\allowbreak div.\allowbreak log.\allowbreak pcl} & stress; ubiquitin or ULP modification & 4 & 2 & 2 \\
\addlinespace[3pt]
545 & \textbf{Ferea 1999}\newline \sourcetext{Ferea TL, Botstein D, Brown PO, Rosenzweig RF. Systematic changes in gene expression patterns following adaptive evolution in yeast. \emph{Proc Natl Acad Sci U S A}} 1999;96(17):9721--6. & \href{https://doi.org/10.1073/pnas.96.17.9721}{doi:10.\allowbreak 1073/\allowbreak pnas.\allowbreak 96.\allowbreak 17.\allowbreak 9721}\newline \href{https://pubmed.ncbi.nlm.nih.gov/10449761/}{PMID 10449761}\newline \href{https://pmc.ncbi.nlm.nih.gov/articles/PMC22277/}{PMC22277} & {\ttfamily\scriptsize 2010.\allowbreak evolution.\allowbreak pcl} & evolution & 4 & -- & 2 \\
\addlinespace[3pt]
551 & \textbf{Keller 2001}\newline \sourcetext{Keller G, Ray E, Brown PO, Winge DR. Haa1, a protein homologous to the copper-regulated transcription factor Ace1, is a novel transcriptional activator. \emph{J Biol Chem}} 2001;276(42):38697--702. & \href{https://doi.org/10.1074/jbc.M107131200}{doi:10.\allowbreak 1074/\allowbreak jbc.\allowbreak M107131200}\newline \href{https://pubmed.ncbi.nlm.nih.gov/11504737/}{PMID 11504737} & {\ttfamily\scriptsize 2010.\allowbreak Keller01.\allowbreak flt.\allowbreak knn.\allowbreak avg.\allowbreak pcl} & metal or metalloid ion stress; starvation; stress; transcription & 4 & -- & 2 \\
\addlinespace[3pt]
554 & \textbf{Miyake 2004}\newline \sourcetext{Miyake T, Reese J, Loch CM, Auble DT, Li R. Genome-wide analysis of ARS (autonomously replicating sequence) binding factor 1 (Abf1p)-mediated transcriptional regulation in Saccharomyces cerevisiae. \emph{J Biol Chem}} 2004;279(33):34865--72. & \href{https://doi.org/10.1074/jbc.M405156200}{doi:10.\allowbreak 1074/\allowbreak jbc.\allowbreak M405156200}\newline \href{https://pubmed.ncbi.nlm.nih.gov/15192094/}{PMID 15192094} & {\ttfamily\scriptsize 2010.\allowbreak Miyake04.\allowbreak flt.\allowbreak knn.\allowbreak avg.\allowbreak pcl} & transcription & 4 & -- & 2 \\
\addlinespace[3pt]
557 & \textbf{Protchenko 2001}\newline \sourcetext{Protchenko O, Ferea T, Rashford J, et al., Philpott CC. Three cell wall mannoproteins facilitate the uptake of iron in Saccharomyces cerevisiae. \emph{J Biol Chem}} 2001;276(52):49244--50. & \href{https://doi.org/10.1074/jbc.M109220200}{doi:10.\allowbreak 1074/\allowbreak jbc.\allowbreak M109220200}\newline \href{https://pubmed.ncbi.nlm.nih.gov/11673473/}{PMID 11673473} & {\ttfamily\scriptsize 2010.\allowbreak Protchenko01.\allowbreak flt.\allowbreak knn.\allowbreak avg.\allowbreak pcl} & metal or metalloid ion stress & 4 & -- & 2 \\
\addlinespace[3pt]
570 & \textbf{Yamamoto 2005}\newline \sourcetext{Yamamoto A, Mizukami Y, Sakurai H. Identification of a novel class of target genes and a novel type of binding sequence of heat shock transcription factor in Saccharomyces cerevisiae. \emph{J Biol Chem}} 2005;280(12):11911--9. & \href{https://doi.org/10.1074/jbc.M411256200}{doi:10.\allowbreak 1074/\allowbreak jbc.\allowbreak M411256200}\newline \href{https://pubmed.ncbi.nlm.nih.gov/15647283/}{PMID 15647283} & {\ttfamily\scriptsize 2010.\allowbreak Yamamoto05.\allowbreak filter.\allowbreak flt.\allowbreak knn.\allowbreak avg.\allowbreak div.\allowbreak log.\allowbreak pcl} & heat shock; transcription & 4 & -- & 2 \\
\addlinespace[3pt]
571 & \textbf{Yun 2000}\newline \sourcetext{Yun CW, Ferea T, Rashford J, et al., Philpott CC. Desferrioxamine-mediated iron uptake in Saccharomyces cerevisiae. Evidence for two pathways of iron uptake. \emph{J Biol Chem}} 2000;275(14):10709--15. & \href{https://doi.org/10.1074/jbc.275.14.10709}{doi:10.\allowbreak 1074/\allowbreak jbc.\allowbreak 275.\allowbreak 14.\allowbreak 10709}\newline \href{https://pubmed.ncbi.nlm.nih.gov/10744769/}{PMID 10744769} & {\ttfamily\scriptsize 2010.\allowbreak Yun00.\allowbreak flt.\allowbreak knn.\allowbreak avg.\allowbreak pcl} & cellular ion homeostasis; metal or metalloid ion stress & 4 & -- & 2 \\
\addlinespace[3pt]
582 & \textbf{Kuhn 2001}\newline \sourcetext{Kuhn KM, DeRisi JL, Brown PO, Sarnow P. Global and specific translational regulation in the genomic response of Saccharomyces cerevisiae to a rapid transfer from a fermentable to a nonfermentable carbon source. \emph{Mol Cell Biol}} 2001;21(3):916--27. & \href{https://doi.org/10.1128/MCB.21.3.916-927.2001}{doi:10.\allowbreak 1128/\allowbreak MCB.\allowbreak 21.\allowbreak 3.\allowbreak 916-\allowbreak 927.\allowbreak 2001}\newline \href{https://pubmed.ncbi.nlm.nih.gov/11154278/}{PMID 11154278}\newline \href{https://pmc.ncbi.nlm.nih.gov/articles/PMC86682/}{PMC86682} & {\ttfamily\scriptsize 2010.\allowbreak Kuhn01.\allowbreak flt.\allowbreak knn.\allowbreak avg.\allowbreak pcl} & carbon utilization; fermentation; respiration & 3 & -- & 2 \\
\addlinespace[3pt]
583 & \textbf{Mizuguchi 2004}\newline \sourcetext{Mizuguchi G, Shen X, Landry J, Wu WH, Sen S, Wu C. ATP-driven exchange of histone H2AZ variant catalyzed by SWR1 chromatin remodeling complex. \emph{Science}} 2004;303(5656):343--8. & \href{https://doi.org/10.1126/science.1090701}{doi:10.\allowbreak 1126/\allowbreak science.\allowbreak 1090701}\newline \href{https://pubmed.ncbi.nlm.nih.gov/14645854/}{PMID 14645854} & {\ttfamily\scriptsize 2010.\allowbreak swr1\_\allowbreak htz1\_\allowbreak ino80.\allowbreak pcl} & chromatin organization & 3 & -- & 2 \\
\addlinespace[3pt]
590 & \textbf{Cullen 2004}\newline \sourcetext{Cullen PJ, Sabbagh W Jr, Graham E, et al., Sprague GF Jr. A signaling mucin at the head of the Cdc42- and MAPK-dependent filamentous growth pathway in yeast. \emph{Genes Dev}} 2004;18(14):1695--708. & \href{https://doi.org/10.1101/gad.1178604}{doi:10.\allowbreak 1101/\allowbreak gad.\allowbreak 1178604}\newline \href{https://pubmed.ncbi.nlm.nih.gov/15256499/}{PMID 15256499}\newline \href{https://pmc.ncbi.nlm.nih.gov/articles/PMC478191/}{PMC478191} & {\ttfamily\scriptsize 2010.\allowbreak glycosylation.\allowbreak pcl} & filamentous growth; protein glycosylation; signaling & 2 & -- & 2 \\
\end{longtable}
\endgroup
\end{landscape}

